\subsection{训练循环}

标准训练循环可以压缩成四步：

\begin{enumerate}
    \item 前向：算预测值和 loss
    \item 反向：`loss.backward()`
    \item 更新：`optimizer.step()`
    \item 清零：\texttt{optimizer.zero\_grad(set\_to\_none=True)}
\end{enumerate}

\begin{lstlisting}
for t in range(num_train_steps):
    x, y = get_batch(B=B)
    pred_y = model(x)
    loss = F.mse_loss(pred_y, y)
    loss.backward()
    optimizer.step()
    optimizer.zero_grad(set_to_none=True)
\end{lstlisting}

\begin{importantbox}{训练循环的本质}
训练循环不是”写起来麻烦”，而是”把前向、反向、更新三种状态转换显式地串起来”。一旦逻辑清楚了，调参和定位 bug 都会容易很多。
\end{importantbox}

